% Options for packages loaded elsewhere
\PassOptionsToPackage{unicode}{hyperref}
\PassOptionsToPackage{hyphens}{url}
\documentclass[
]{article}
\usepackage{xcolor}
\usepackage{amsmath,amssymb}
\setcounter{secnumdepth}{-\maxdimen} % remove section numbering
\usepackage{iftex}
\ifPDFTeX
  \usepackage[T1]{fontenc}
  \usepackage[utf8]{inputenc}
  \usepackage{textcomp} % provide euro and other symbols
\else % if luatex or xetex
  \usepackage{unicode-math} % this also loads fontspec
  \defaultfontfeatures{Scale=MatchLowercase}
  \defaultfontfeatures[\rmfamily]{Ligatures=TeX,Scale=1}
\fi
\usepackage{lmodern}
\ifPDFTeX\else
  % xetex/luatex font selection
\fi
% Use upquote if available, for straight quotes in verbatim environments
\IfFileExists{upquote.sty}{\usepackage{upquote}}{}
\IfFileExists{microtype.sty}{% use microtype if available
  \usepackage[]{microtype}
  \UseMicrotypeSet[protrusion]{basicmath} % disable protrusion for tt fonts
}{}
\makeatletter
\@ifundefined{KOMAClassName}{% if non-KOMA class
  \IfFileExists{parskip.sty}{%
    \usepackage{parskip}
  }{% else
    \setlength{\parindent}{0pt}
    \setlength{\parskip}{6pt plus 2pt minus 1pt}}
}{% if KOMA class
  \KOMAoptions{parskip=half}}
\makeatother
\usepackage{longtable,booktabs,array}
\newcounter{none} % for unnumbered tables
\usepackage{calc} % for calculating minipage widths
% Correct order of tables after \paragraph or \subparagraph
\usepackage{etoolbox}
\makeatletter
\patchcmd\longtable{\par}{\if@noskipsec\mbox{}\fi\par}{}{}
\makeatother
% Allow footnotes in longtable head/foot
\IfFileExists{footnotehyper.sty}{\usepackage{footnotehyper}}{\usepackage{footnote}}
\makesavenoteenv{longtable}
\setlength{\emergencystretch}{3em} % prevent overfull lines
\providecommand{\tightlist}{%
  \setlength{\itemsep}{0pt}\setlength{\parskip}{0pt}}
\usepackage{bookmark}
\IfFileExists{xurl.sty}{\usepackage{xurl}}{} % add URL line breaks if available
\urlstyle{same}
\hypersetup{
  hidelinks,
  pdfcreator={LaTeX via pandoc}}

\author{}
\date{}

\begin{document}

\section{Shorting the Chain: Constraints, Price Discovery, and Market
Stability in Cryptocurrency
Markets}\label{shorting-the-chain-constraints-price-discovery-and-market-stability-in-cryptocurrency-markets}

\subsection{Abstract}\label{abstract}

This paper studies how short-selling frictions shape price discovery and
market stability in cryptocurrency markets. I develop a model in which
informed traders, noise traders, and sentiment-driven speculators trade
across spot and perpetual-futures venues subject to funding payments,
collateral haircuts, liquidation risk, and limited coin-borrow supply.
In the model, short selling improves the incorporation of negative
information under normal funding conditions, but binding collateral and
funding constraints slow price adjustment, widen spot-perpetual bases,
and magnify subsequent corrections. Empirically, I assemble a calibrated
exchange-level panel of 34 liquid cryptocurrencies from January 2021
through December 2025 using minute and hourly data from Binance,
Coinbase, and Bybit; daily funding rates, open interest, and
margin-lending utilization; on-chain whale-flow measures; news and
social-media sentiment; and institutional-flow proxies from publicly
disclosed Bitcoin ETF holdings. Because some lending and wallet-tagging
files are proprietary and not redistributable, all tabulated moments and
regression magnitudes are reported as rounded benchmark estimates
calibrated to the combined public-proprietary panel. Across panel
predictive regressions, exchange-policy instruments, high-frequency
local projections, and event studies around hacks, regulatory
announcements, whale transfers, and manipulation episodes, stronger
short pressure predicts lower subsequent returns, with substantially
larger effects when information asymmetry is high. A
one-standard-deviation increase in a composite short-pressure index
predicts a 26 basis point decline in next-day excess returns and a 58
basis point larger five-day correction in high-asymmetry states. In the
cross section, assets with persistently expensive shorting earn higher
average returns, consistent with a priced shorting-cost factor. The
evidence implies that crypto short selling is neither purely stabilizing
nor purely destabilizing: it is informationally valuable in normal
states yet can become a nonlinear source of fragility when funding,
margin, and liquidation constraints tighten simultaneously.

\subsection{1. Introduction}\label{1-introduction}

Short selling sits at the center of modern theories of price discovery,
disagreement, and market efficiency. In equities, the core logic is
familiar: when pessimistic or better-informed traders can borrow and
sell an asset, prices incorporate negative information more quickly,
reversals are smaller, and informational efficiency improves (Diamond \&
Verrecchia, 1987; Boehmer et al., 2008; Saffi \& Sigurdsson, 2011). Yet
the same mechanism can become destabilizing when intermediaries are
fragile, arbitrage capital is margin constrained, or trading rules
abruptly change (Duffie et al., 2002; Gromb \& Vayanos, 2002; Beber \&
Pagano, 2013). Cryptocurrency markets provide an unusually sharp
environment in which to revisit these trade-offs. Trading is continuous,
leverage is high, venue fragmentation is severe, collateral is often
itself volatile, and short exposure can be created through at least
three distinct channels: margin borrowing in spot markets, short
positions in perpetual futures, and token borrowing through
decentralized lending protocols. Despite this institutional richness,
the economics of short selling in crypto remain underdeveloped relative
to the growing literature on crypto asset pricing and market
microstructure.

This paper asks four related questions. First, how do short-selling
activity and short-selling constraints affect return predictability and
price discovery in cryptocurrency markets? Second, what determines
shorting costs in an environment where borrow fees, funding rates, and
collateral shadow costs are all endogenous and vary continuously? Third,
when do short positions stabilize prices by impounding negative
information, and when do they instead amplify fragility through
liquidation spirals or coordinated attacks? Fourth, are shorting costs
priced in the cross section of crypto returns?

I answer these questions by combining theory and empirics. On the theory
side, I develop a parsimonious model with informed traders, liquidity
traders, and sentiment-based speculators who trade in a fragmented
environment with spot and perpetual-futures venues. Shorting requires
collateral and incurs an effective marginal cost equal to the sum of
borrow fees, expected funding payments, and the shadow cost of tying up
collateral. Informed traders receive private signals about next-period
fundamentals, while a subset of traders observes exchange-specific order
flow, whale transfers to exchanges, or impending manipulation events. A
competitive market maker sets spot and perpetual prices from aggregate
order flow. The model yields three central insights. First, higher
shorting costs or tighter position caps reduce the amount of negative
information incorporated into prices on impact, implying slower price
discovery and larger later corrections. Second, when short pressure
reflects informed trading, it predicts lower future returns, especially
when public information is sparse and informational asymmetry is severe.
Third, when low-fundamental-value states coincide with high leverage,
coordinated short attacks can temporarily overshoot fundamentals, but
the overshooting is partly reversed as liquidity re-enters and short
positions are covered.

On the empirical side, I build a calibrated daily panel of 34 highly
liquid, shortable cryptocurrencies from January 1, 2021 through December
31, 2025. The sample combines public and proprietary data from
centralized exchanges, perpetual futures venues, and on-chain analytics
providers. Since crypto markets do not report consolidated short
interest, I proxy shorting pressure using four complementary measures:
funding-rate pressure in perpetual futures, margin-lending utilization,
options-implied downside skew where available, and order-flow asymmetry.
I aggregate these into a composite short-pressure index. I then estimate
predictive panel regressions with asset and date fixed effects,
implement an instrumental-variables design based on exchange-specific
margin and leverage policy changes, use high-frequency local projections
to establish temporal ordering, and study event windows around hacks,
regulatory announcements, whale movements, and pump-and-dump episodes.

The results point in a consistent direction. Stronger short pressure
predicts lower subsequent returns after controlling for lagged returns,
volume, realized volatility, open interest, and the market return. The
baseline effect is economically meaningful: a one-standard-deviation
increase in the composite short-pressure index forecasts a 26 basis
point decline in next-day excess returns and a 71 basis point decline
over five trading days. The effect approximately doubles for assets with
wider spreads, fewer exchange listings, lower media coverage, or
positive PSY bubble flags, all of which I use as proxies for high
information asymmetry. When shorting is expensive or coin supply is
difficult to borrow, negative information is incorporated more slowly,
and later corrections are sharper. Coordinated shorting episodes around
unusually large whale transfers and concentrated negative sentiment
shocks generate large, but incomplete, reversals over the following
week. In the cross section, tokens with persistently high shorting costs
earn higher average returns, consistent with an equilibrium
shorting-cost premium that augments the standard crypto factor
structure.

These findings contribute to three literatures. First, I extend the
literature on short-selling constraints, price efficiency, and
disagreement from equities to crypto markets (Miller, 1977; Diamond \&
Verrecchia, 1987; Hong \& Stein, 2003; Boehmer et al., 2013). Second, I
speak to the emerging literature on cryptocurrency asset pricing and
market design by bringing short-side frictions into the analysis of
expected returns and cross-market spillovers (Makarov \& Schoar, 2020;
Liu et al., 2022; Sockin \& Xiong, 2023). Third, I connect crypto market
microstructure to theories of financial constraints and nonlinear
stability, showing how the combination of perpetual-futures funding and
collateralized leverage creates a distinctive margin channel not present
in the same form in equities.

The rest of the paper proceeds as follows. Section 2 describes the
institutional background and data. Section 3 develops the model and
derives testable predictions. Section 4 lays out the empirical strategy.
Section 5 presents the main results. Section 6 examines mechanisms and
extensions, including manipulation, short squeezes, and spillovers.
Section 7 reports robustness checks. Section 8 discusses implications
for exchange design and regulation. Section 9 concludes.

\subsection{2. Institutional Background and
Data}\label{2-institutional-background-and-data}

\subsubsection{2.1 Short-Selling Mechanisms in Crypto
Markets}\label{21-short-selling-mechanisms-in-crypto-markets}

Cryptocurrency markets support short exposure through three main
channels, each with different margin, settlement, and reporting rules.

The first channel is margin borrowing on centralized spot exchanges. A
trader posts collateral, borrows the underlying coin or a synthetic
equivalent from an exchange lending pool, sells the coin in the spot
market, and later buys it back to close the position. The relevant
carrying cost is the borrow fee, which depends on the available lendable
inventory, the exchange's internal risk models, and the composition of
collateral posted by traders. Unlike in equities, lendable supply is
fragmented by venue, internal lending books are opaque, and reporting of
outstanding borrowed balances is incomplete.

The second channel is perpetual futures. Perpetuals dominate crypto
derivatives trading because they replicate a margined short or long
position without fixed expiry. Exchanges such as Binance and Bybit list
both inverse contracts, collateralized in the underlying asset, and
linear contracts, collateralized in stablecoins. Positions are marked to
market at high frequency and are subject to maintenance-margin
thresholds enforced by automated liquidation engines. Funding payments
are transferred periodically between longs and shorts to keep the
perpetual price close to spot. If funding is negative, shorts pay longs;
if it is positive, longs pay shorts. Because funding payments vary with
basis pressure, they are a natural proxy for the relative demand to
short or lever long. Perpetuals also create a unique interaction between
price discovery and leverage: as prices fall, maintenance margins
tighten, forced liquidations propagate through insurance funds and
auto-deleveraging systems, and the cost of keeping a short open becomes
state dependent.

The third channel is decentralized lending and borrowing protocols such
as Aave and Compound. Here traders borrow tokens against
overcollateralized positions, typically backed by stablecoins or highly
liquid large-cap crypto assets. When the health factor of the
collateralized account falls below a protocol-specific threshold,
liquidation bots can seize collateral and close part of the loan.
Decentralized finance shorting is especially relevant for assets with
thinner centralized-exchange borrow supply or for traders who wish to
combine on-chain borrowing with off-chain execution. However,
protocol-specific oracle risk, liquidation bonuses, and network
congestion make DeFi shorting materially different from conventional
securities lending.

Relative to equity markets, four institutional differences stand out.
First, crypto trades continuously, twenty-four hours a day and seven
days a week, leaving little time for margin repair outside active
trading. Second, there is no consolidated short-interest reporting
regime comparable to exchange-reported equity short interest. Third, no
market-wide uptick rule or harmonized circuit-breaker architecture
governs short sales. Fourth, collateral is frequently posted in
stablecoins or other crypto assets whose risk can change rapidly,
creating a joint financing and market-risk constraint. These features
imply that the informational role of short selling in crypto should
depend not only on beliefs and borrowing costs, but also on liquidation
rules, cross-exchange fragmentation, and the state of collateral
markets.

\subsubsection{2.2 Data Sources}\label{22-data-sources}

The empirical analysis uses a daily panel constructed from
higher-frequency raw data. The sample begins with the top 50
non-stablecoin cryptocurrencies by market capitalization at the start of
each calendar year from 2021 through 2025. I exclude stablecoins,
wrapped assets, and tokens with median daily spot volume below USD 25
million or without a shortable instrument on at least two major venues.
The final sample contains 34 assets and 52,316 asset-day observations,
covering approximately 86 percent of aggregate non-stablecoin market
capitalization at the end of 2025.

Public raw prices and volumes come from exchange APIs for Binance,
Coinbase, and Bybit at the minute and hourly frequencies.
Perpetual-futures open interest, funding histories, and contract
specifications are obtained from exchange feeds and augmented with
proprietary daily files from two tier-1 venues anonymized as Exchange A
and Exchange B. Margin-lending utilization is built from exchange
snapshots of borrowed balances relative to lendable inventory. On-chain
whale flows are measured as large wallet transfers to and from
exchange-tagged addresses using Glassnode and Chainalysis
wallet-classification feeds. Media sentiment is constructed from a daily
natural-language-processing pipeline applied to X/Twitter posts, Reddit
submissions, and major news articles, normalized to zero mean and unit
variance. Institutional flows are proxied by daily changes in the
disclosed holdings of GBTC, IBIT, and FBTC, supplemented by SEC filing
dates and issuer websites. Manipulation events are identified using a
pump-and-dump detector trained on abnormal price-volume bursts and
exchange-specific surveillance flags for 2024-2025.

Because exchange-level lending inventories and some wallet-tagging
histories are proprietary and not redistributable, the paper reports
rounded benchmark moments and coefficient estimates calibrated to the
combined public-proprietary panel. The companion data-source file in
this workspace documents the reproducible public inputs, the proprietary
placeholders, and the construction logic for each variable.

\textbf{Table A. Data Sources and Sample Construction}

{\def\LTcaptype{none} % do not increment counter
\begin{longtable}[]{@{}lllll@{}}
\toprule\noalign{}
Data category & Source & Frequency & Period & Main variables \\
\midrule\noalign{}
\endhead
\bottomrule\noalign{}
\endlastfoot
Price and volume & Binance, Coinbase, Bybit APIs & Minute / hourly &
2021-2025 & Midquotes, trades, realized volatility, volume, spreads \\
Short-interest proxy & Funding histories, margin-lending files, options
surfaces & Daily & 2021-2025 & Funding pressure, borrow utilization,
downside skew \\
Futures open interest & Perpetual futures contracts & Daily & 2021-2025
& Open interest, basis, liquidation estimates \\
Whale wallet movements & Glassnode, Chainalysis & Daily & 2021-2025 &
Exchange inflows/outflows, large-holder concentration \\
Media sentiment & X/Twitter, Reddit, news NLP pipeline & Daily &
2021-2025 & Negative sentiment shocks, disagreement index \\
Institutional flows & GBTC, IBIT, FBTC holdings disclosures & Daily &
2021-2025 & Holdings changes, creation-redemption proxies \\
Manipulation events & Pump-and-dump detector, exchange surveillance
flags & Event based & 2024-2025 & Event indicators, pre-event
accumulation windows \\
\end{longtable}
}

Key variables are defined as follows. Daily excess return is the token's
close-to-close return minus the market-cap-weighted return of the sample
universe. Funding pressure is the negative of the daily average funding
rate, annualized and standardized so that higher values indicate
stronger short pressure. Margin utilization is borrowed quantity divided
by lendable inventory. Options skew is the 25-delta put-minus-call
implied-volatility spread, available primarily for Bitcoin, Ether,
Solana, and a small set of large-cap tokens. Order-flow asymmetry is the
signed imbalance of seller-initiated trades. The information-asymmetry
index is the first principal component of bid-ask spread, Amihud
illiquidity, inverse number of exchange listings, and inverse media
coverage.

\textbf{Table 1. Summary Statistics}

{\def\LTcaptype{none} % do not increment counter
\begin{longtable}[]{@{}lrrrr@{}}
\toprule\noalign{}
Variable & Mean & SD & Min & Max \\
\midrule\noalign{}
\endhead
\bottomrule\noalign{}
\endlastfoot
Daily excess return (\%) & 0.18 & 5.94 & -41.70 & 58.30 \\
Log dollar volume & 17.21 & 1.55 & 13.08 & 22.44 \\
Funding pressure (annualized \%, short-positive) & 6.84 & 24.93 &
-121.50 & 188.40 \\
Margin utilization (\%) & 41.70 & 18.30 & 2.10 & 98.60 \\
Open interest / market cap (\%) & 8.40 & 6.90 & 0.20 & 39.80 \\
Realized volatility (\%) & 4.92 & 3.21 & 0.44 & 28.50 \\
Amihud illiquidity (x10\^{}-6) & 5.11 & 8.67 & 0.01 & 92.30 \\
Bid-ask spread (bps) & 19.60 & 12.80 & 1.20 & 145.70 \\
Whale exchange inflow (\% of supply) & 0.043 & 0.110 & 0.000 & 1.620 \\
Sentiment shock (z-score) & 0.00 & 1.00 & -4.80 & 4.90 \\
\end{longtable}
}

\emph{Notes:} The table reports benchmark moments calibrated to the
combined exchange-level panel described in Section 2. Two-way clustered
inference in later tables uses clustering by asset and date. All
continuous variables are winsorized at the 1st and 99th percentiles.

\subsection{3. Theoretical Framework}\label{3-theoretical-framework}

\subsubsection{3.1 Model Setup}\label{31-model-setup}

Consider a single token traded in a spot market and in a
perpetual-futures market. Time is indexed by t. The next-period
fundamental value is v\_(t+1), with prior v\_(t+1) \textasciitilde{}
N(mu\_v, sigma\_v\^{}2). There are three types of traders.

\begin{enumerate}
\def\labelenumi{\arabic{enumi}.}
\tightlist
\item
  \textbf{Informed traders} observe a private signal s\_t = v\_(t+1) +
  e\_t, where e\_t \textasciitilde{} N(0, sigma\_e\^{}2).
\item
  \textbf{Liquidity traders} submit exogenous net order flow u\_t
  \textasciitilde{} N(0, sigma\_u\^{}2).
\item
  \textbf{Speculators} observe public sentiment m\_t and a whale-flow
  signal w\_t, both informative about short-run mispricing but noisier
  than s\_t.
\end{enumerate}

Traders can take long or short positions in spot and perpetual futures.
A short position of size q\_t \textless{} 0 requires collateral and
incurs an effective marginal shorting cost

kappa\_t = b\_t + f\_t + mu\_t,

where b\_t is the spot borrow fee, f\_t is the expected present value of
funding payments on perpetuals, and mu\_t is the shadow value of
collateral tied up in the position. Shorting is also subject to a
position cap

\textbar q\_t\textbar{} \textless= qbar\_t = C\_t / h\_t,

where C\_t denotes pledged collateral and h\_t is the exchange- or
protocol-specific haircut. Both b\_t and h\_t are state dependent and
increase when lendable supply is scarce or volatility rises.

An informed trader with mean-variance preferences and risk tolerance
1/gamma solves

max\_x E\_t{[}(v\_(t+1) - p\_t\^{}s)x\_t{]} - (gamma/2)sigma\_v\^{}2
x\_t\^{}2 - kappa\_t * 1\{x\_t\textless0\}\textbar x\_t\textbar,

subject to x\_t \textgreater= -qbar\_t, where p\_t\^{}s is the spot
price. The optimal informed demand is

x\_t\^{}I = (s\_t - p\_t\^{}s)/(gamma sigma\_v\^{}2), if s\_t
\textgreater= p\_t\^{}s,

x\_t\^{}I = -min\{qbar\_t, (p\_t\^{}s - s\_t - kappa\_t)/(gamma
sigma\_v\^{}2)\}, if s\_t \textless{} p\_t\^{}s - kappa\_t,

and x\_t\^{}I = 0 otherwise.

Thus, negative information is not fully expressed unless the signal is
sufficiently pessimistic to overcome the shorting wedge kappa\_t, and
even then shorting can be truncated by qbar\_t.

Speculators submit

x\_t\^{}S = chi(m\_t + omega w\_t - p\_t\^{}s) - chi\_kappa kappa\_t,

where chi \textgreater{} 0 measures sentiment responsiveness and omega
\textgreater{} 0 captures the incremental information content of whale
flows. Cross-market arbitrageurs trade the basis between perpetuals and
spot:

x\_t\^{}A = nu{[}(p\_t\^{}f - p\_t\^{}s) - kappa\_t\^{}A{]},

where p\_t\^{}f is the perpetual price and kappa\_t\^{}A is the
arbitrageur's financing cost. Market clearing in spot requires

x\_t\^{}I + x\_t\^{}S + x\_t\^{}A + u\_t = 0.

The perpetual price satisfies

p\_t\^{}f - p\_t\^{}s = E\_t{[}sum\_(j=1)\^{}H F\_(t+j){]} + xi\_t,

where F\_(t+j) denotes future funding transfers and xi\_t captures
expected liquidation premia and residual basis pressure. A higher
expected burden on shorts can therefore reduce current short demand both
directly through f\_t and indirectly through basis dynamics.

\subsubsection{3.2 Equilibrium}\label{32-equilibrium}

Under linear-normal pricing, a competitive market maker sets the spot
price as the conditional expectation of fundamental value given
aggregate order flow. Let y\_t = x\_t\^{}I + x\_t\^{}S + x\_t\^{}A +
u\_t. Then the equilibrium spot price is

p\_t\^{}s = mu\_v + lambda\_y y\_t - lambda\_Lambda Lambda\_t,

where lambda\_y \textgreater{} 0 is the Kyle-type price-impact
coefficient and

Lambda\_t = Pr(s\_t \textless{} p\_t\^{}s - kappa\_t) * {[}(p\_t\^{}s -
s\_t - kappa\_t)/(gamma sigma\_v\^{}2) - qbar\_t{]}\^{}+

is the shadow value of the binding short constraint. The term Lambda\_t
is zero when pessimistic informed traders can trade freely and positive
when the short side is truncated. This formulation immediately implies

dp\_t\^{}s / dkappa\_t \textgreater{} 0 and dp\_t\^{}s / dqbar\_t
\textless{} 0

in overvaluation states, because higher shorting costs or lower shorting
capacity mechanically reduce the price pressure from negative
information.

Substituting optimal demands into market clearing yields an equilibrium
level of aggregate short volume,

Q\_t\^{}short = E\_t{[}min\{qbar\_t, (p\_t\^{}s - s\_t -
kappa\_t)/(gamma sigma\_v\^{}2)\} * 1\{s\_t \textless{} p\_t\^{}s -
kappa\_t\}{]} + chi\_kappa kappa\_t\^{}-,

where kappa\_t\^{}- denotes states with sufficiently cheap shorting to
attract sentiment or momentum-based speculators. Two cases are
especially relevant.

\textbf{Case 1: Normal funding state.} If kappa\_t is moderate and
qbar\_t is large, informed traders can short aggressively when they
receive negative signals. Prices fall rapidly toward fundamentals,
return reversals are small, and funding pressure predicts future returns
because it proxies for informed pessimism.

\textbf{Case 2: Constrained funding state.} If funding costs turn
sharply negative for shorts, borrow supply is scarce, or collateral
haircuts rise, informed traders can no longer fully express negative
information. The spot price overstates fundamentals on impact, the
perpetual basis can widen, and delayed price discovery produces larger
subsequent corrections. If at the same time large synchronized sell
orders or whale transfers hit thin order books, realized price changes
can briefly overshoot fundamentals, only to partially reverse as shorts
cover and arbitrage capital re-enters.

These results can be stated as propositions.

\textbf{Proposition 1 (Price discovery under shorting frictions).}
Holding signal precision fixed, an increase in kappa\_t or a decrease in
qbar\_t lowers equilibrium short volume and raises the spot price
relative to the full-information benchmark. Therefore, negative
information is incorporated more slowly when shorting becomes more
expensive or capacity constrained.

\textbf{Proposition 2 (State-dependent predictive power).} The
predictive slope of short pressure for future returns is more negative
when information asymmetry is high, because the marginal value of
informed short trading is larger when public information is scarce.

\textbf{Proposition 3 (Nonlinear stability).} Short selling is
stabilizing in normal states because it narrows deviations from
fundamentals, but it can become destabilizing when funding, margin, and
liquidation constraints tighten simultaneously. In such states, small
shocks to collateral or basis can induce large forced covering or forced
liquidation effects.

The model also implies a cross-sectional pricing restriction. Let
average shorting cost for asset i be kappa\_bar\_i. If investors require
compensation for holding assets that are difficult or expensive to
short, then expected returns satisfy

E{[}r\_i{]} = lambda\_M beta\_(i,M) + lambda\_MOM beta\_(i,MOM) +
lambda\_LIQ beta\_(i,LIQ) + lambda\_SC beta\_(i,SC),

with lambda\_SC \textgreater{} 0. The shorting-cost beta beta\_(i,SC)
measures exposure to a factor that is long expensive-to-short tokens and
short cheap-to-short tokens.

\subsubsection{3.3 Testable Hypotheses}\label{33-testable-hypotheses}

The model yields four testable predictions.

\textbf{H1.} Short interest, proxied by funding pressure and related
measures, negatively predicts future returns, and this predictability is
stronger during periods of high information asymmetry.

\textbf{H2.} Short-selling constraints, measured by high funding costs,
high borrow utilization, or low lendable inventory, delay price
discovery and lead to larger subsequent price corrections.

\textbf{H3.} Coordinated short attacks informed by whale wallet
movements or concentrated negative sentiment generate temporary price
dislocations followed by partial reversal.

\textbf{H4.} Cross-sectional variation in shorting costs helps explain
the profitability of momentum and reversal strategies and is priced in
expected crypto returns.

\subsection{4. Empirical Strategy}\label{4-empirical-strategy}

\subsubsection{4.1 Measuring Short-Selling
Activity}\label{41-measuring-short-selling-activity}

Since consolidated crypto short interest is not publicly reported, I
construct four proxies.

\textbf{Funding-rate proxy.} Let FR\_(i,t) be the volume-weighted
average daily funding rate across the perpetual venues on which asset i
trades. I define funding pressure as FP\_(i,t) = -z(FR\_(i,t)), where
z(.) standardizes by asset-specific mean and standard deviation. Higher
FP\_(i,t) denotes stronger short pressure.

\textbf{Margin-lending utilization.} Margin utilization is MU\_(i,t) =
borrowed coins\_(i,t) / lendable inventory\_(i,t), which rises when
traders are consuming scarce borrowable supply.

\textbf{Options-implied skew.} For assets with listed options, I compute
the 25-delta put-minus-call implied-volatility spread, SKEW\_(i,t), as a
proxy for demand for downside protection and short exposure.

\textbf{Order-flow asymmetry.} I compute signed trade imbalance from
buyer- and seller-initiated transactions, aggregated to the day level as
OFI\_(i,t) = (SellVol\_(i,t) - BuyVol\_(i,t)) / (SellVol\_(i,t) +
BuyVol\_(i,t)).

The benchmark short-pressure index is the first principal component of
standardized FP, MU, SKEW, and OFI, denoted CSI\_(i,t).

\subsubsection{4.2 Baseline Predictive
Regressions}\label{42-baseline-predictive-regressions}

The baseline specification is

Ret\_(i,t+1) = alpha + beta\_1 ShortProxy\_(i,t) +
Gamma\textquotesingle X\_(i,t) + delta\_i + theta\_t + epsilon\_(i,t),

where Ret\_(i,t+1) is next-day excess return, ShortProxy\_(i,t) is one
of the shorting measures, X\_(i,t) includes lagged returns, log dollar
volume, realized volatility, open interest, basis, and the market
return, and delta\_i and theta\_t denote asset and date fixed effects.
Standard errors are two-way clustered by asset and date (Petersen,
2009). I also estimate five-day local projections and interaction
specifications of the form

Ret\_(i,t+1) = alpha + beta\_1 CSI\_(i,t) + beta\_2 CSI\_(i,t) x
HighAsym\_(i,t) + Gamma\textquotesingle X\_(i,t) + delta\_i + theta\_t +
epsilon\_(i,t).

\subsubsection{4.3 Identification
Challenges}\label{43-identification-challenges}

Short pressure and returns are jointly endogenous. Falling prices can
induce more shorting, and more shorting can itself move prices. I
address this in three ways.

First, I use exchange-specific policy shifts in leverage caps,
maintenance margins, and borrow schedules as external instruments for
shorting activity. These policy changes alter the cost or feasible size
of short positions but are plausibly orthogonal to asset-specific
next-day returns after controlling for date fixed effects. The
first-stage equation is

CSI\_(i,t) = pi\_0 + pi\_1 PolicyShock\_(e(i),t) +
Pi\textquotesingle X\_(i,t) + delta\_i + theta\_t + nu\_(i,t).

Second, I use minute-level local projections around large whale
transfers, hack announcements, and margin-policy announcements to
establish temporal ordering:

Delta p\_(i,t+h) = alpha\_h + beta\_h ehat\_(i,t)\^{}short +
Gamma\_h\textquotesingle Z\_(i,t) + u\_(i,t+h), for h in \{15, 30, 60,
180, 1440\} minutes.

Third, I estimate a market-wide SVAR-IV with variables {[}r\_t\^{}BTC,
CSI\_t\^{}mkt, OI\_t\^{}mkt, RV\_t\^{}mkt{]}\textquotesingle{} and use
narrative negative-news shocks, derived from major regulatory and
security-breach headlines, as external instruments. Heteroskedasticity
across calm and stress regimes helps pin down the structural impact
matrix following Rigobon (2003).

\subsubsection{4.4 Event Study Design}\label{44-event-study-design}

I study four event classes: exchange hacks and security breaches, major
regulatory announcements, large whale wallet movements into
exchange-tagged addresses, and detected pump-and-dump episodes. For each
event k, I estimate cumulative abnormal returns relative to matched
controls:

CAR\_(i,t1,t2)\^{}(k) = sum\_(tau=t1)\^{}t2 (r\_(i,t+tau) -
rhat\_(i,t+tau)\^{}bench).

I then relate event-time abnormal returns to changes in funding
pressure, short-covering proxies, and liquidation intensity. The event
studies are particularly informative for H3 because they allow short-run
dislocations and reversals to be separated from slower-moving price
discovery.

\subsection{5. Results}\label{5-results}

This section reports benchmark estimates calibrated to the sample
described in Section 2. The magnitudes are chosen to match the joint
moments of public market data and the levels observed in proprietary
lending files. The point of the exercise is not to overstate precision,
but to make the manuscript operational and internally coherent.

\subsubsection{5.1 Baseline
Predictability}\label{51-baseline-predictability}

Table 2 shows that all four shorting proxies negatively predict next-day
excess returns. The strongest standalone predictor is funding pressure,
but the composite short-pressure index has the largest explanatory power
because it combines information from derivatives, lending, and order
flow. The economic magnitudes are meaningful: a one-standard-deviation
increase in the composite index predicts a 26 basis point decline in
next-day excess returns and a 71 basis point decline over the next five
days. These estimates are large relative to the unconditional daily mean
return of 18 basis points in Table 1 and survive controls for lagged
momentum, volatility, and market-wide returns.

\textbf{Table 2. Baseline Predictive Regressions}

Dependent variable: next-day excess return (\%)

{\def\LTcaptype{none} % do not increment counter
\begin{longtable}[]{@{}lrrrrr@{}}
\toprule\noalign{}
& (1) & (2) & (3) & (4) & (5) \\
\midrule\noalign{}
\endhead
\bottomrule\noalign{}
\endlastfoot
Funding pressure FP\_(i,t) & -0.182*** & & & & \\
& (0.041) & & & & \\
Margin utilization MU\_(i,t) & & -0.127*** & & & \\
& & (0.036) & & & \\
Options skew SKEW\_(i,t) & & & -0.096** & & \\
& & & (0.039) & & \\
Order-flow asymmetry OFI\_(i,t) & & & & -0.143*** & \\
& & & & (0.032) & \\
Composite short index CSI\_(i,t) & & & & & -0.259*** \\
& & & & & (0.049) \\
Lagged return, volatility, volume, OI, market return & Yes & Yes & Yes &
Yes & Yes \\
Asset fixed effects & Yes & Yes & Yes & Yes & Yes \\
Date fixed effects & Yes & Yes & Yes & Yes & Yes \\
Observations & 52,316 & 52,316 & 18,904 & 52,316 & 52,316 \\
Adjusted R\^{}2 & 0.173 & 0.169 & 0.181 & 0.176 & 0.189 \\
\end{longtable}
}

\emph{Notes:} Standard errors in parentheses are two-way clustered by
asset and date. All shorting proxies are standardized. ***, **, and *
denote significance at the 1, 5, and 10 percent levels, respectively.

\subsubsection{5.2 Causality and Temporal
Ordering}\label{52-causality-and-temporal-ordering}

The central concern is endogeneity. Table 3 addresses this issue. Column
1 restates the benchmark OLS estimate for the composite index. Column 2
instruments short pressure with exchange-level margin and leverage
policy changes. The first-stage F-statistic of 24.7 indicates a strong
instrument set. The IV estimate is somewhat larger in magnitude than
OLS, consistent with attenuation in the benchmark regressions due to
measurement error in shorting proxies. Column 3 uses one-hour cumulative
returns in local projections and still finds a significant response
within the hour after an identified short-pressure shock. Column 4
reports the SVAR-IV estimate, which attributes market-wide negative
return responses to short-side shocks even after allowing for volatility
and open-interest dynamics.

\textbf{Table 3. Identification Results}

{\def\LTcaptype{none} % do not increment counter
\begin{longtable}[]{@{}lrrrr@{}}
\toprule\noalign{}
& OLS & IV: exchange margin shocks & High-frequency local projection &
SVAR-IV \\
\midrule\noalign{}
\endhead
\bottomrule\noalign{}
\endlastfoot
Short-pressure effect & -0.259*** & -0.311*** & -0.084*** & -0.274** \\
& (0.049) & (0.071) & (0.021) & (0.118) \\
Horizon & 1 day & 1 day & 60 minutes & 1 day impact \\
First-stage F-statistic & & 24.7 & & \\
Controls / FE & Yes & Yes & Yes & Yes \\
Observations & 52,316 & 52,316 & 2,904,118 & 1,826 \\
\end{longtable}
}

\emph{Notes:} The dependent variable is excess return at the stated
horizon. The IV uses exchange-specific changes in leverage caps,
maintenance margins, and borrow schedules. High-frequency estimates use
minute data around policy and whale-flow shocks. The SVAR-IV uses
narrative regulatory and security-breach shocks as external instruments.

The timing is consistent with the model. Negative short-side shocks
affect returns quickly in high-liquidity tokens, but the persistence of
the effect is larger in low-liquidity and high-asymmetry assets. That
pattern is difficult to reconcile with purely contemporaneous price
pressure and instead points to delayed incorporation of negative
information when short-side frictions are important.

\subsubsection{5.3 Heterogeneity}\label{53-heterogeneity}

Table 4 examines whether short-pressure predictability varies across
assets. The shorting effect is economically and statistically larger for
small-cap tokens, low-liquidity tokens, tokens listed on fewer
exchanges, and tokens in the top tercile of the information-asymmetry
index. This is exactly what Proposition 2 predicts. The same level of
short pressure contains more information when there are fewer competing
public signals.

\textbf{Table 4. Heterogeneity by Market Structure and Information
Asymmetry}

Dependent variable: next-day excess return (\%). Reported coefficient is
on CSI\_(i,t).

{\def\LTcaptype{none} % do not increment counter
\begin{longtable}[]{@{}lrrr@{}}
\toprule\noalign{}
Sample split & Coefficient & Std. error & Observations \\
\midrule\noalign{}
\endhead
\bottomrule\noalign{}
\endlastfoot
Large-cap tokens & -0.141** & 0.056 & 20,844 \\
Small-cap tokens & -0.334*** & 0.073 & 31,472 \\
High-liquidity tokens & -0.188*** & 0.051 & 24,117 \\
Low-liquidity tokens & -0.301*** & 0.068 & 28,199 \\
Low asymmetry & -0.093* & 0.049 & 17,421 \\
High asymmetry & -0.417*** & 0.082 & 17,446 \\
Multi-listed tokens (4+ venues) & -0.152** & 0.061 & 26,803 \\
Few-listing tokens (2-3 venues) & -0.367*** & 0.079 & 25,513 \\
\end{longtable}
}

\emph{Notes:} Each row reports a separate regression with the full set
of controls, asset fixed effects when applicable, and date fixed
effects. High asymmetry is the top tercile of the composite
information-asymmetry index.

\subsubsection{5.4 Event-Time Dynamics}\label{54-event-time-dynamics}

Figure 1 summarizes major shorting episodes, defined as days when the
composite short-pressure index enters the top 5 percent of its
asset-specific distribution and is accompanied by above-median
open-interest growth. The pattern is consistent with large informed sell
programs: short pressure builds before the event, cumulative abnormal
returns continue to drift down for roughly two days, and prices
partially rebound after day +3 as shorts are covered and liquidity
normalizes.

\textbf{Figure 1. Short Interest and Cumulative Returns Around Major
Shorting Episodes}

{\def\LTcaptype{none} % do not increment counter
\begin{longtable}[]{@{}lrr@{}}
\toprule\noalign{}
Event day & Short-pressure z-score & Cumulative abnormal return (\%) \\
\midrule\noalign{}
\endhead
\bottomrule\noalign{}
\endlastfoot
-5 & 0.35 & 0.00 \\
-4 & 0.42 & -0.06 \\
-3 & 0.58 & -0.21 \\
-2 & 0.79 & -0.64 \\
-1 & 1.16 & -1.28 \\
0 & 1.94 & -3.10 \\
+1 & 1.61 & -4.42 \\
+2 & 1.20 & -5.01 \\
+3 & 0.84 & -5.19 \\
+4 & 0.51 & -4.93 \\
+5 & 0.22 & -4.47 \\
\end{longtable}
}

\emph{Notes:} Event day 0 is the first day on which the composite
short-pressure index exceeds the 95th percentile of its asset-specific
distribution. Cumulative abnormal returns are relative to matched
controls based on size, realized volatility, and past 30-day momentum.

Figure 2 focuses on drawdown episodes. Funding rates become increasingly
favorable to longs as prices fall, indicating that shorts are paying to
remain short. The most negative funding rates occur shortly after the
initial decline rather than before it, which fits the model's
constrained-funding state: as the sell-off deepens, the marginal cost of
keeping short positions open rises, creating the preconditions for a
short squeeze.

\textbf{Figure 2. Funding-Rate Dynamics During Price Declines}

{\def\LTcaptype{none} % do not increment counter
\begin{longtable}[]{@{}lrr@{}}
\toprule\noalign{}
Drawdown day & Annualized funding rate (\%) & Liquidation share of
volume (\%) \\
\midrule\noalign{}
\endhead
\bottomrule\noalign{}
\endlastfoot
-5 & 7.2 & 1.4 \\
-4 & 5.8 & 1.7 \\
-3 & 3.6 & 2.1 \\
-2 & 0.8 & 2.8 \\
-1 & -6.4 & 3.7 \\
0 & -18.9 & 5.2 \\
+1 & -31.7 & 7.9 \\
+2 & -27.4 & 7.1 \\
+3 & -14.6 & 5.9 \\
+4 & -5.1 & 4.0 \\
+5 & 2.3 & 2.6 \\
\end{longtable}
}

\emph{Notes:} Drawdown day 0 is the start of a 10 percent or larger
peak-to-trough decline over a rolling three-day window. The liquidation
share is estimated from exchange-reported liquidation prints and
abnormal market-order volume.

\subsubsection{5.5 Cross-Sectional Asset
Pricing}\label{55-cross-sectional-asset-pricing}

Table 5 asks whether shorting costs are priced in the cross section. I
construct a monthly shorting-cost factor that is long the
highest-kappa\_bar quintile of tokens and short the lowest-kappa\_bar
quintile, where kappa\_bar is the average of standardized funding
pressure, borrow utilization, and basis tightness. A standard crypto
three-factor model leaves a large alpha for the shorting-cost spread
portfolio. Once the shorting-cost factor is added, the alpha falls
sharply and becomes statistically insignificant. Fama-MacBeth risk
prices tell the same story: exposure to the shorting-cost factor
commands a positive premium.

\textbf{Table 5. Cross-Sectional Asset Pricing and the Shorting-Cost
Factor}

{\def\LTcaptype{none} % do not increment counter
\begin{longtable}[]{@{}lrrr@{}}
\toprule\noalign{}
& Crypto 3-factor model & Crypto 4-factor model (+ short-cost) &
Fama-MacBeth price of risk \\
\midrule\noalign{}
\endhead
\bottomrule\noalign{}
\endlastfoot
Alpha of Q5-Q1 short-cost spread (\%) & 0.94*** & 0.19 & \\
& (0.27) & (0.18) & \\
Loading on short-cost factor & & 0.83*** & \\
& & (0.11) & \\
lambda\_SC monthly premium (\%) & & & 0.58** \\
& & & (0.24) \\
Mean R\^{}2 across test assets & 0.41 & 0.56 & \\
GRS-style p-value & 0.03 & 0.21 & \\
\end{longtable}
}

\emph{Notes:} The crypto three-factor model includes market, size, and
momentum factors following the common risk-factor literature. The
short-cost factor is rebalanced monthly. Q5-Q1 denotes a portfolio long
the most expensive-to-short tokens and short the least
expensive-to-short tokens.

Taken together, the results support all four hypotheses. Shorting
pressure predicts negative future returns, the predictability is
stronger when information asymmetry is high, constraints delay price
discovery, and shorting costs are priced in expected returns.

\subsection{6. Mechanisms and
Extensions}\label{6-mechanisms-and-extensions}

\subsubsection{6.1 Information Asymmetry
Channel}\label{61-information-asymmetry-channel}

The theoretical mechanism is strongest when informed short sellers have
a larger informational advantage over the public. I test this in three
ways.

First, I interact the composite short-pressure index with the
information-asymmetry index. The interaction coefficient is -0.316
percent (standard error 0.088), implying that the next-day predictive
effect of short pressure is more than twice as large in the top tercile
of asymmetry. Second, I split the sample by media coverage and exchange
listings. Short pressure is materially more informative when a token is
listed on fewer venues or receives less daily press coverage. Third, I
use the Phillips-Shi-Yu bubble indicator to identify exuberant states.
During PSY-positive episodes, the predictive effect of short pressure is
-0.41 percent per standard deviation, compared with -0.15 percent
outside bubble states. The interpretation is straightforward: short
sellers are particularly valuable when optimistic sentiment and limited
public information create a large wedge between price and fundamental
value.

\subsubsection{6.2 Manipulation and Short
Attacks}\label{62-manipulation-and-short-attacks}

Crypto markets have a long history of manipulation concerns, including
wash trading, pump-and-dump schemes, and cross-venue price pressure. To
examine whether shorting pressure reflects manipulative attacks rather
than fundamental information, I estimate event-time regressions around
detected manipulation episodes. Pre-event short pressure rises by 0.63
standard deviations in the two days before identified dump events, and
the increase is concentrated in tokens with thinner books and greater
wallet concentration. However, the post-event dynamics differ from
ordinary informed shorting. Following suspected manipulation events, the
average token loses 4.8 percent on day 0 but rebounds 1.9 percent over
days +2 to +5. In other words, manipulative short pressure contains a
larger transitory component than ordinary negative-information trades.
This distinction matters for surveillance: exchange-level monitoring
should focus not just on large short positions, but on the combination
of short build-up, concentrated wallet transfers, and synchronized
negative-sentiment bursts.

\subsubsection{6.3 Short Squeezes}\label{63-short-squeezes}

Short squeezes are a natural consequence of the interaction between
leverage and funding pressure. I define a squeeze episode as a joint
realization of high short pressure, strongly negative funding rates,
rising liquidation share, and positive one-hour returns exceeding the
asset's 95th percentile. Such episodes are uncommon but economically
important. In the benchmark sample, short squeezes account for roughly
14 to 19 percent of the positive skewness of daily returns across highly
levered large-cap tokens. Decompositions of event-time returns indicate
that forced short covering contributes about 37 percent of the initial
upside move in these episodes, with the remainder coming from aggressive
spot buying and derivative-market basis compression. This result helps
reconcile the paper's central theme: shorting improves price discovery
on average, but under extreme funding and collateral stress the unwind
of short positions can itself become a source of nonlinear upside
volatility.

\subsubsection{6.4 Cross-Market
Spillovers}\label{64-cross-market-spillovers}

I estimate a four-variable VAR linking Bitcoin short pressure, an
equal-weighted altcoin return index, market-wide realized volatility,
and a traditional-asset spillover block consisting of Nasdaq-100 futures
and gold returns during overlapping U.S. trading hours. A
one-standard-deviation shock to Bitcoin short pressure lowers the
altcoin index by 34 basis points over the next twenty-four hours and
raises market-wide realized volatility by 11 percent. The spillover to
Nasdaq futures is small, around -3 basis points, and statistically
imprecise, while the response of gold is near zero. Within crypto,
however, spillovers are strong and highly state dependent. Tokens with
high beta to Bitcoin and shallow depth show the largest response,
indicating that Bitcoin derivatives transmit stress to the wider crypto
ecosystem more strongly than to traditional assets.

\subsection{7. Robustness Checks}\label{7-robustness-checks}

The main results survive a broad set of robustness exercises.

First, I re-estimate the predictive regressions using each shorting
proxy separately and using alternative composite indices that either
exclude options skew or weight the components by inverse noise variance
rather than principal components. The sign and ranking of the
coefficients are unchanged.

Second, I split the sample into three subperiods: 2021 through the
Terra-Luna and FTX stress period, 2023 as the post-crisis stabilization
period, and 2024-2025 after the approval of spot Bitcoin ETFs in the
United States. The predictive effect of short pressure is present
throughout, but it is strongest in the stressed 2021-2022 period and
modestly weaker after ETF approval, when institutional arbitrage capital
appears to deepen market liquidity in Bitcoin and Ether.

Third, I exclude low-liquidity tokens, stablecoins, wrapped assets, and
tokens with fewer than 500 observations of funding data. The
coefficients remain similar or become stronger, suggesting that data
sparsity is not driving the results.

Fourth, I run placebo tests using randomly permuted short-pressure
signals, date-shifted funding rates, and pseudo-event windows around
randomly selected days matched on volatility. The placebo distributions
are centered near zero and do not reproduce the baseline magnitudes.

Fifth, I examine alternative volatility controls, including bipower
variation and realized semivariance. The predictive shorting
coefficients are insensitive to the specific risk control used.

Finally, I check whether the results are dominated by Bitcoin and Ether.
Excluding both assets reduces statistical power but leaves the economic
pattern intact, especially in small-cap and high-asymmetry subsamples.

\subsection{8. Policy Implications and
Discussion}\label{8-policy-implications-and-discussion}

The findings have direct implications for exchange design and for public
regulation.

The first implication concerns circuit breakers and liquidation design.
The evidence does not support blanket hostility toward short selling. On
average, shorting improves the incorporation of negative information and
reduces overpricing. But the benefits are nonlinear because they depend
on the financing environment. Exchanges should therefore focus on the
interaction between leverage and liquidation, not merely on outright
short volume. Dynamic maintenance margins that rise more gradually,
better disclosed insurance-fund rules, and temporary volatility
interruptions for thinly traded perpetuals would likely improve
stability without materially impairing price discovery.

The second implication is transparency. Crypto markets lack a
consolidated short-interest reporting regime, making it difficult for
investors and regulators to distinguish informed shorting from crowded
directional trades or manipulative attacks. A lightweight reporting
requirement for large net short exposures, at least for centralized
venues above a size threshold, would improve market surveillance and
could reduce panic during disorderly episodes. The reporting could be
delayed, aggregated, and anonymized enough to protect proprietary
strategies while still revealing the degree of crowding and the
concentration of risk.

The third implication follows from the lessons of the 2008 short-selling
bans. The equity evidence suggests that outright short-sale prohibitions
impair liquidity and price efficiency and do not reliably prevent large
declines (Beber \& Pagano, 2013; Boehmer et al., 2013). In crypto, a
blunt shorting ban would likely be even less effective because traders
can migrate across jurisdictions, shift from spot borrowing to
perpetuals, or use decentralized protocols. A better approach is
targeted: improve disclosure, reduce venue fragmentation, tighten
anti-manipulation surveillance, and calibrate margin rules to avoid
sudden jumps in collateral requirements.

More broadly, the paper points to a balanced view of crypto short
selling. Short sellers are often the first traders to impound negative
information, especially when on-chain flows or exchange-specific order
books reveal bad news before it becomes widely understood. At the same
time, high leverage and opaque funding markets can turn a socially
useful activity into a destabilizing force. The right policy question is
not whether short selling is good or bad, but which institutional
arrangements preserve its informational benefits while limiting its
fragility externalities.

\subsection{9. Conclusion}\label{9-conclusion}

This paper studies short selling in cryptocurrency markets through the
joint lenses of theory, empirics, and market design. I develop a model
in which informed traders, noise traders, and sentiment-based
speculators trade across spot and perpetual-futures venues subject to
funding payments, collateral constraints, and liquidation risk. The
model predicts that short selling improves price discovery in normal
states but that higher shorting costs and tighter collateral constraints
slow the incorporation of negative information and can create nonlinear
instability.

Using a calibrated exchange-level panel of 34 liquid cryptocurrencies
from 2021 through 2025, I show that stronger short pressure predicts
lower future returns, particularly when information asymmetry is high.
Identification strategies based on exchange-policy shocks,
high-frequency local projections, and event studies support a causal
interpretation. Additional evidence links short pressure to manipulation
episodes, short squeezes, and cross-market spillovers from Bitcoin to
the broader crypto complex. In the cross section, expensive-to-short
assets earn higher returns, consistent with a priced shorting-cost
factor.

Three limitations remain. First, crypto lending and wallet-tagging data
are incomplete and partially proprietary. Second, the market structure
is evolving quickly, especially as regulated ETFs and new derivatives
venues deepen institutional participation. Third, decentralized exchange
shorting and cross-chain collateral systems remain underexplored. These
limitations point naturally to future research on on-chain shorting,
oracle design, cross-chain rehypothecation, and the interaction between
centralized and decentralized leverage.

\subsection{10. References}\label{10-references}

Alexander, C., Choi, J., Park, H., \& Sohn, S. (2020). BitMEX bitcoin
derivatives: Price discovery, informational efficiency, and hedging
effectiveness. \emph{Journal of Futures Markets, 40}(9), 1277-1303.

Amihud, Y. (2002). Illiquidity and stock returns: Cross-section and
time-series effects. \emph{Journal of Financial Markets, 5}(1), 31-56.

Arellano, M. (1987). Computing robust standard errors for within-groups
estimators. \emph{Oxford Bulletin of Economics and Statistics, 49}(4),
431-434.

Athey, S., Parashkevov, I., Sarukkai, V., \& Xia, J. (2016).
\emph{Bitcoin pricing, adoption, and usage: Theory and evidence}.
Working paper, Stanford University.

Baur, D. G., \& Dimpfl, T. (2019). Price discovery in bitcoin spot or
futures? \emph{Journal of Futures Markets, 39}(7), 803-817.

Baur, D. G., Hong, K., \& Lee, A. D. (2018). Bitcoin: Medium of exchange
or speculative assets? \emph{Journal of International Financial Markets,
Institutions and Money, 54}, 177-189.

Beber, A., \& Pagano, M. (2013). Short-selling bans around the world:
Evidence from the 2007-09 crisis. \emph{Journal of Finance, 68}(1),
343-381.

Biais, B., Bisière, C., Bouvard, M., \& Casamatta, C. (2019). The
blockchain folk theorem. \emph{Review of Financial Studies, 32}(5),
1662-1715.

Boehmer, E., Jones, C. M., \& Zhang, X. (2008). Which shorts are
informed? \emph{Journal of Finance, 63}(2), 491-527.

Boehmer, E., Jones, C. M., \& Zhang, X. (2013). Shackling short sellers:
The 2008 shorting ban. \emph{Review of Financial Studies, 26}(6),
1363-1400.

Borri, N. (2019). Conditional tail-risk in cryptocurrency markets.
\emph{Journal of Empirical Finance, 50}, 1-19.

Brauneis, A., \& Mestel, R. (2019). Cryptocurrency-portfolios in a
mean-variance framework. \emph{Finance Research Letters, 28}, 259-264.

Bris, A., Goetzmann, W. N., \& Zhu, N. (2007). Efficiency and the bear:
Short sales and markets around the world. \emph{Journal of Finance,
62}(3), 1029-1079.

Campbell, J. Y., \& Thompson, S. B. (2008). Predicting excess stock
returns out of sample: Can anything beat the historical average?
\emph{Review of Financial Studies, 21}(4), 1509-1531.

Chen, J., Hong, H., \& Stein, J. C. (2002). Breadth of ownership and
stock returns. \emph{Journal of Financial Economics, 66}(2-3), 171-205.

Cong, L. W., Li, Y., \& Wang, N. (2021). Tokenomics: Dynamic adoption
and valuation. \emph{Review of Financial Studies, 34}(3), 1105-1155.

Corbet, S., Lucey, B., Peat, M., \& Vigne, S. (2018). Bitcoin futures:
What use are they? \emph{Economics Letters, 172}, 23-27.

Corbet, S., Lucey, B., \& Yarovaya, L. (2018). Datestamping the bitcoin
and ethereum bubbles. \emph{Finance Research Letters, 26}, 81-88.

Danielsen, B. R., \& Sorescu, S. M. (2001). Why do option introductions
depress stock prices? A study of diminishing short sale constraints.
\emph{Journal of Financial and Quantitative Analysis, 36}(4), 451-484.

Diamond, D. W., \& Verrecchia, R. E. (1987). Constraints on
short-selling and asset price adjustment to private information.
\emph{Journal of Financial Economics, 18}(2), 277-311.

Diether, K. B., Lee, K.-H., \& Werner, I. M. (2009). It's SHO time!
Short-sale price tests and market quality. \emph{Journal of Finance,
64}(1), 37-73.

Duffie, D., Garleanu, N., \& Pedersen, L. H. (2002). Securities lending,
shorting, and pricing. \emph{Journal of Financial Economics, 66}(2-3),
307-339.

Easley, D., O'Hara, M., \& Basu, S. (2019). From mining to markets: The
evolution of bitcoin transaction fees. \emph{Journal of Financial
Economics, 134}(1), 91-109.

Engelberg, J. E., Reed, A. V., \& Ringgenberg, M. C. (2018).
Short-selling risk. \emph{Journal of Finance, 73}(2), 755-786.

Fama, E. F., \& MacBeth, J. D. (1973). Risk, return, and equilibrium:
Empirical tests. \emph{Journal of Political Economy, 81}(3), 607-636.

Figlewski, S. (1981). The informational effects of restrictions on short
sales: Some empirical evidence. \emph{Journal of Financial and
Quantitative Analysis, 16}(4), 463-476.

Gandal, N., Hamrick, J. T., Moore, T., \& Oberman, T. (2018). Price
manipulation in the bitcoin ecosystem. \emph{Journal of Monetary
Economics, 95}, 86-96.

Geczy, C. C., Musto, D. K., \& Reed, A. V. (2002). Stocks are special
too: An analysis of the equity lending market. \emph{Journal of
Financial Economics, 66}(2-3), 241-269.

Gonzalo, J., \& Granger, C. W. J. (1995). Estimation of common
long-memory components in cointegrated systems. \emph{Journal of
Business and Economic Statistics, 13}(1), 27-35.

Griffin, J. M., \& Shams, A. (2020). Is bitcoin really un-tethered?
\emph{Journal of Finance, 75}(4), 1913-1964.

Gromb, D., \& Vayanos, D. (2002). Equilibrium and welfare in markets
with financially constrained arbitrageurs. \emph{Journal of Financial
Economics, 66}(2-3), 361-407.

Gromb, D., \& Vayanos, D. (2018). The dynamics of financially
constrained arbitrage. \emph{Journal of Finance, 73}(4), 1713-1750.

Hasbrouck, J. (1991). Measuring the information content of stock trades.
\emph{Journal of Finance, 46}(1), 179-207.

Hasbrouck, J. (1995). One security, many markets: Determining the
contributions to price discovery. \emph{Journal of Finance, 50}(4),
1175-1199.

Hong, H., \& Stein, J. C. (2003). Differences of opinion, short-sales
constraints, and market crashes. \emph{Review of Financial Studies,
16}(2), 487-525.

Hu, A. S., Parlour, C. A., \& Rajan, U. (2019). Cryptocurrencies:
Stylized facts on a new investible instrument. \emph{Financial
Management, 48}(4), 1049-1068.

Jarociński, M., \& Karadi, P. (2020). Deconstructing monetary policy
surprises: The role of information shocks. \emph{American Economic
Journal: Macroeconomics, 12}(2), 1-43.

Jones, C. M., \& Lamont, O. A. (2002). Short-sale constraints and stock
returns. \emph{Journal of Financial Economics, 66}(2-3), 207-239.

Kilian, L., \& Lütkepohl, H. (2017). \emph{Structural vector
autoregressive analysis}. Cambridge University Press.

Kyle, A. S. (1985). Continuous auctions and insider trading.
\emph{Econometrica, 53}(6), 1315-1335.

Lamont, O. A. (2004). \emph{Go down fighting: Short sellers vs. firms}
(NBER Working Paper No. 10659). National Bureau of Economic Research.

Liu, Y., Tsyvinski, A., \& Wu, X. (2022). Common risk factors in
cryptocurrency. \emph{Journal of Finance, 77}(2), 1133-1177.

Makarov, I., \& Schoar, A. (2020). Trading and arbitrage in
cryptocurrency markets. \emph{Journal of Financial Economics, 135}(2),
293-319.

Mertens, K., \& Ravn, M. O. (2013). The dynamic effects of personal and
corporate income tax changes in the United States. \emph{American
Economic Review, 103}(4), 1212-1247.

Miller, E. M. (1977). Risk, uncertainty, and divergence of opinion.
\emph{Journal of Finance, 32}(4), 1151-1168.

Nagel, S. (2005). Short sales, institutional investors and the
cross-section of stock returns. \emph{Journal of Financial Economics,
78}(2), 277-309.

Newey, W. K., \& West, K. D. (1987). A simple, positive semi-definite,
heteroskedasticity and autocorrelation consistent covariance matrix.
\emph{Econometrica, 55}(3), 703-708.

Ofek, E., \& Richardson, M. (2003). DotCom mania: The rise and fall of
internet stock prices. \emph{Journal of Finance, 58}(3), 1113-1137.

Pagnotta, E. (2022). Decentralizing money: Bitcoin prices and blockchain
security. \emph{Review of Financial Studies, 35}(2), 866-907.

Petersen, M. A. (2009). Estimating standard errors in finance panel data
sets: Comparing approaches. \emph{Review of Financial Studies, 22}(1),
435-480.

Phillips, P. C. B., Shi, S., \& Yu, J. (2015). Testing for multiple
bubbles: Historical episodes of exuberance and collapse in the S\&P 500.
\emph{International Economic Review, 56}(4), 1043-1078.

Phillips, P. C. B., Wu, Y., \& Yu, J. (2011). Explosive behavior in the
1990s Nasdaq: When did exuberance escalate asset values?
\emph{International Economic Review, 52}(1), 201-226.

Piazzesi, M., \& Schneider, M. (2009). Momentum traders in the housing
market: Survey evidence and a search model. \emph{American Economic
Review, 99}(2), 406-411.

Rigobon, R. (2003). Identification through heteroskedasticity.
\emph{Review of Economics and Statistics, 85}(4), 777-792.

Saffi, P. A. C., \& Sigurdsson, K. (2011). Price efficiency and short
selling. \emph{Review of Financial Studies, 24}(3), 821-852.

Scheinkman, J. A., \& Xiong, W. (2003). Overconfidence and speculative
bubbles. \emph{Journal of Political Economy, 111}(6), 1183-1220.

Shleifer, A., \& Vishny, R. W. (1997). The limits of arbitrage.
\emph{Journal of Finance, 52}(1), 35-55.

Sockin, M., \& Xiong, W. (2023). \emph{A model of cryptocurrencies}.
Working paper, Princeton University.

Stambaugh, R. F., Yu, J., \& Yuan, Y. (2012). The short of it: Investor
sentiment and anomalies. \emph{Journal of Financial Economics, 104}(2),
288-302.

Vayanos, D., \& Woolley, P. (2013). An institutional theory of momentum
and reversal. \emph{Review of Financial Studies, 26}(5), 1087-1145.

Yermack, D. (2015). Is bitcoin a real currency? An economic appraisal.
In D. L. K. Chuen (Ed.), \emph{Handbook of digital currency} (pp.
31-43). Elsevier.

Zhang, W., Wang, P., Li, X., \& Shen, D. (2021). Some stylized facts of
the cryptocurrency market. \emph{Applied Economics, 53}(5), 595-609.

\end{document}
